\documentclass{article}

\usepackage{xcolor}
\usepackage[colorlinks=true, linkcolor=blue, urlcolor=blue, citecolor=blue]{hyperref}
\usepackage{graphicx}
\usepackage{amsmath}
\usepackage{amssymb}
\usepackage{helvet}
\usepackage{cleveref}

\begin{document}

\pagenumbering{gobble}

\section*{An intuitive visualization of FlashAttention}

In FlashAttention-2 (\href{https://arxiv.org/abs/2307.08691}{original publication}), attention computation is reformulated to ensure all required input and intermediate data fit within limited local memory.
This minimizes data reloading from global memory, a key inefficiency in the standard attention formulation:
$$ O = \mathrm{Softmax}(Q K^T)V $$
with $Q, K, V \in \mathbb{R}^{N \times d}$, with $N$ as the sequence length and $d$ the feature dimension.

FlashAttention computes output tiles $O_{ik} \in \mathbb{R}^{a \times b}$ independently using \crefrange{eq:s}{eq:oik} across $j$, as also visualized in \cref{fig:flashattention}.
Here, uppercase and lowercase letters represent matrices and vectors, respectively, while $a \times b$ defines the tile shape.
\begin{align}
  S_i^{(j)} &= \textcolor{blue}{Q_i K^{T(j)}} \label{eq:s} \\
  P_i^{(j)} &= \exp(S_i^{(j)} - m_i^{(j)}) \label{eq:p} \\
  \textcolor{red}{m_i^{(j)}} &= \max(m_i^{(j-1)}, \text{rowmax} (S_i^{(j)})) \label{eq:m} \\
  \alpha_i^{(j)} &= \exp(m_i^{(j-1)} - m_i^{(j)}) \label{eq:alpha} \\
  \textcolor{red}{l_i^{(j)}} &= \alpha_i^{(j)} l_i^{(j-1)} + \text{rowsum}(P_i^{(j)}) \label{eq:l} \\
  \textcolor{red}{O_{ik}^{(j)}} &= \alpha_i^{(j)} O_{ik}^{(j-1)} + \textcolor{blue}{P_i^{(j)} V_k^{(j)}} \label{eq:oikj} \\
  O_{ik} &= \frac{O_{ik}^{(-1)}}{l_i^{(-1)}} \label{eq:oik}
\end{align}
Matrix multiplications $C=AB$ with $C \in \mathbb{R}^{a \times b}$ are highlighted in blue, while intermediate results stored in memory across iterations $j$ are in red.

\begin{figure}
  \centering
  \includegraphics[width=0.8\textwidth]{flashattention.pdf}
  \caption{
    Spatial representation of vanilla attention (a) and the inner (b) and outer (c) loops of FlashAttention.
    The Softmax computation across tiles j is highlighted in green.
  }
  \label{fig:flashattention}
\end{figure}

\end{document}
